\section{问题定义与任务形式化}
\subsection{任务背景}
本文以临床长文本表型识别任务作为主要研究载体。该任务的目标是：给定一份病历文本，预测其中是否包含若干具有临床意义的表型标签。与单标签分类相比，这一任务更接近真实医疗场景，因为同一名患者可能同时表现出多个风险因素、慢性疾病或精神行为特征，模型必须在同一份文本中识别多种可能并存的标签。

在本文所使用的任务设定中，输出空间由十四个表型标签组成，包括高级癌症、高级心脏病、高级肺病、酒精滥用、慢性神经退行性疾病、慢性疼痛、痴呆、抑郁、发育迟缓、依从性差、肥胖、其他药物滥用、精神分裂谱系障碍以及不确定标签。这一标签集合既包含明确的医学概念，也包含行为和精神状态相关标签，因此在语义上具有较强异质性。这种异质性恰好为可信输出层研究提供了合适场景：不同标签的证据形式不同，不同样本的噪声水平也并不一致。

\subsection{形式化表示}
设输入临床文本为 $x_i$，其对应标签向量为 $\mathbf{y}_i \in \{0,1\}^{K}$，其中 $K=14$ 表示标签维数。长上下文编码器 $E(\cdot)$ 将输入文本映射到 $d$ 维特征空间：
\[
\mathbf{z}_i = E(x_i), \qquad \mathbf{z}_i \in \mathbb{R}^{d}.
\]
将全部样本特征堆叠为矩阵 $\mathbf{X} \in \mathbb{R}^{N\times d}$，标签堆叠为矩阵 $\mathbf{Y} \in \mathbb{R}^{N\times K}$。传统方法通常学习一个确定性映射 $f(\mathbf{z}_i)$ 输出 logits 或类别分数，而本文进一步定义可信输出映射：
\[
F(\mathbf{z}_i)=\big(\boldsymbol{\mu}_i,\log \boldsymbol{\sigma}_i^2,\mathbf{w}_i\big),
\]
其中 $\boldsymbol{\mu}_i$ 表示预测均值，$\boldsymbol{\sigma}_i^2$ 表示输入相关方差，$\mathbf{w}_i$ 为由方差诱导的样本权重，用于后续结构化求解。

这一形式化定义体现了本文与普通多标签分类的本质差异。标准模型通常只回答“输出是什么”；而本文的方法同时回答“输出有多不确定”和“这种不确定性是否应当反过来影响最后一层的估计方式”。这使得模型具备了从判别走向可信判别的可能性。

\subsection{为何需要可信输出层}
在临床文本任务中，所有样本等权的假设通常并不合理。部分病历文本记录充分、证据直接、标签边界清晰；另一些病历则存在记录缺失、叙事跳跃、否定与既往史混杂或标签边界模糊等问题。如果输出层仍然把所有样本视为同等可靠，那么训练过程容易受高噪声样本干扰，最终导致模型在困难样本上过度自信，在简单样本上又缺乏区分能力。

由此，本文提出如下核心判断：对于临床长文本任务，仅靠更强的编码器并不足以解决可信性问题；在编码器之后构建能够感知样本噪声水平并据此调整输出函数的可信输出层，是一条独立且必要的研究路径。本文的主要工作正是围绕这一判断展开。

\subsection{研究假设}
为使问题具备可实现性，本文在方法上采用以下三个假设。第一，\encoder{} 已经提供了较强的长文本表示能力，因此在当前研究阶段，可以先冻结编码器，把研究重点集中在输出层。第二，临床样本之间存在显著异方差，不同样本及不同标签维度的噪声水平并不一致，因此需要显式建模输入相关方差。第三，最后一层参数若仅由梯度下降隐式学习，难以充分利用样本权重所携带的统计信息；引入 Ridge/WLS 一类结构化求解器，有助于增强数值稳定性和判别稳健性。

这些假设共同决定了本文的研究边界：本文不是重新设计医学基础模型，也不是把整个系统做成全贝叶斯模型，而是在“强编码器 + 可信输出层”这一中间层次上寻找可实现、可解释、可扩展的解决方案。

\subsection{研究边界}
本文的研究边界主要体现在三个方面。第一，本文默认长上下文编码器已经能够提供较稳定的临床语义表示，因此重点讨论的是输出层可信建模，而不是重新训练更大的基础模型。第二，本文当前关注的是特征级多标签预测场景，相关分析首先围绕 Phenotype 任务展开，尚未扩展到所有临床 NLP 子任务。第三，本文主要建模输入相关的偶然不确定性，并利用该信息改进最后一层参数估计，尚未把认知不确定性校准、拒识机制和完整贝叶斯化分析纳入统一框架。

明确这些边界，有助于更准确地理解本文贡献所在：本文不是试图一次性解决临床可信预测中的全部问题，而是针对“强编码器之后如何构建更可靠输出层”这一核心问题给出可实现的答案。
